\chapter{工质建模方法}

\section{联合循环机组工质建模对象与总体思路}

联合循环机组由气体循环和蒸汽循环耦合构成，机组内部同时存在燃料气、空气、燃烧产物烟气、水以及水蒸气等多类工质。不同工质在热力性质、组分、状态确定方式以及能量转换方面存在明显差异。燃气轮机侧主要涉及气体混合物的压缩、燃烧和膨胀过程，余热锅炉和汽水系统侧则主要涉及水和水蒸气的加压、加热、蒸发、过热、膨胀和凝结过程。因此，在建立联合循环机组能量分析与火用分析模型之前，有必要首先建立统一、清晰且可复用的工质流股建模方法，为后续压气机、燃烧室、燃气透平、余热锅炉、汽轮机和凝汽器等设备模型提供统一的工质输入和输出接口。

本文的工质建模方法以流股为基本对象，首先需要确定各类工质的状态参数。对于水和水蒸气工质，其状态通常由两个独立状态参数确定。常用输入组合包括压力和温度 $(P,T)$、压力和干度 $(P,x)$、压力和比焓 $(P,h)$ 以及压力和比熵 $(P,s)$。对于气体工质，在上述两个独立状态参数的基础上还需要考虑气体组分。因此，工质状态建模的核心是根据不同输入组合确定其余参数，为后续分析提供基础。

联合循环机组性能分析涉及多个设备和多种工质。为了保证不同设备模型之间能够稳定传递状态信息，工质状态模型需要满足统一性、可扩展性等要求。因此分别采用本文采用 IAPWS-IF97 物性库计算水和水蒸气的热力性质，采用 CoolProp 开源物性库辅助计算气体工质的热力性质。在气体工质物性计算中，本文以标准状态作为基准状态，即
\[
T_{\mathrm{ref}}=273.15\,\mathrm{K}, \qquad
P_{\mathrm{ref}}=101325\,\mathrm{Pa}.
\]
在该基准状态下定义混合气的基准比焓和基准比熵。

在流股的封装方面，本文在程序实现中将工质模块划分为统一流股状态、气体工质、水和水蒸气工质、参考环境以及实测数据流股构造等部分。其核心思想是：先由基础物性模型计算单个工质状态，再将状态封装为统一的流股对象，最后将流股对象作为设备模型的输入和输出。通过该方法，可以使联合循环机组中不同设备之间的工质状态传递具有统一的数据结构和接口形式，从而提高模型的可读性、可维护性和工程适用性，为后续的自动化建模和数据驱动的建模提供了坚实的基础。

\section{流股状态的统一描述方法}

在联合循环机组中，工质在不同设备之间以流股形式传递。无论是进入压气机的空气、进入燃烧室的燃料气、流经燃气透平和余热锅炉的烟气，还是汽水系统中的给水、饱和水、饱和蒸汽和过热蒸汽，本质上都可以看作携带质量、能量和火用的流动介质。为了避免不同设备模型之间重复定义状态参数，本文采用统一流股对象对各类工质状态进行描述。

统一流股对象的核心思想是：将工质的质量流量、温度、压力、比焓、比熵、组成、能量流率和火用流率等参数封装为同一类状态对象，并将其作为设备模型的输入和输出。这样可以使压气机、燃烧室、燃气透平、余热锅炉、汽轮机和凝汽器等设备模型采用一致的接口形式，从而提高模型的可读性和可扩展性。

对于任一流股 $k$，其状态可表示为
\[
\mathcal{F}_k
=
\left\{
\dot{m}_k,\,
T_k,\,
P_k,\,
h_k,\,
s_k,\,
\boldsymbol{x}_k,\,
e_{x,k},\,
\right\},
\]
其中，$\dot{m}_k$ 为质量流量，$T_k$ 为温度，$P_k$ 为压力，$h_k$ 为比焓，$s_k$ 为比熵，$\boldsymbol{x}_k$ 为组分，$e_{x,k}$ 为比火用。其中，比焓 $h_k$ 反映单位质量工质所携带的热力学能量，比熵 $s_k$ 则用于描述工质状态变化的不可逆性和火用水平。基于此可以计算比火用 $e_{x,k}$。若以参考环境为基准，比火用可表示为
\[
e_{x,k}
=
\left(h_k-h_{0,k}\right)
-
T_0\left(s_k-s_{0,k}\right),
\]
其中，$T_0$ 为环境参考温度，$h_{0,k}$ 和 $s_{0,k}$ 分别为第 $k$ 条流股对应工质在参考环境下的比焓和比熵。

联合循环机组的性能分析通常包括能量分析和火用分析两个层次。能量分析以热力学第一定律为基础，重点关注能量在不同设备之间的传递与转换；火用分析以热力学第二定律为基础，重点关注能量品质和不可逆损失。因此，在流股状态建模中，需要同时建立能量流和火用流的统一表达。

对于任一流股，其焓流率可表示为
\[
\dot{H}
=
\dot{m}h.
\]

对于具有多个入口和出口的稳态设备，忽略动能和位能变化时，其能量平衡方程可写为
\[
\dot{Q}
-
\dot{W}
+
\sum_{\mathrm{in}}\dot{m}_{\mathrm{in}}h_{\mathrm{in}}
-
\sum_{\mathrm{out}}\dot{m}_{\mathrm{out}}h_{\mathrm{out}}
=
0,
\]
其中，$\dot{Q}$ 为设备与外界的换热量，$\dot{W}$ 为设备对外输出的功。若设备为压气机或泵，则通常表现为外界对设备输入功；若设备为燃气透平或汽轮机，则表现为设备对外输出功。

火用流反映的是流股相对于参考环境所具有的最大理论做功能力。对于不考虑动能和位能的流股，其物理比火用为
\[
e_x
=
\left(h-h_0\right)
-
T_0\left(s-s_0\right).
\]
对应的火用流率为
\[
\dot{E}_x
=
\dot{m}
\left[
\left(h-h_0\right)
-
T_0\left(s-s_0\right)
\right].
\]

若进一步考虑流股动能和位能，则总比火用可写为
\[
e_{x,\mathrm{tot}}
=
\left(h-h_0\right)
-
T_0\left(s-s_0\right)
+
\frac{c^2}{2}
+
gz,
\]
其中，$c$ 为流速，$g$ 为重力加速度，$z$ 为高度。对于本文所研究的联合循环机组热力系统，设备进出口高度差和流速变化对整体能量分析影响较小，因此可以忽略不计。

对于稳态开口系统，火用平衡方程可表示为
\[
\sum_{\mathrm{in}}\dot{E}_{x,\mathrm{in}}
-
\sum_{\mathrm{out}}\dot{E}_{x,\mathrm{out}}
+
\sum_j
\left(
1-\frac{T_0}{T_j}
\right)
\dot{Q}_j
-
\dot{W}
=
\dot{E}_{x,\mathrm{dest}},
\]
其中，$T_j$ 为第 $j$ 个换热边界温度，$\dot{Q}_j$ 为通过该边界传递的热量，$\dot{E}_{x,\mathrm{dest}}$ 为火用损失率。根据 Gouy--Stodola 定理，火用损失与熵产之间满足
\[
\dot{E}_{x,\mathrm{dest}}
=
T_0\dot{S}_{\mathrm{gen}},
\]
其中，$\dot{S}_{\mathrm{gen}}$ 为设备内部不可逆过程产生的熵产率。

对于绝热设备，若忽略与外界的换热，则火用平衡可简化为
\[
\sum_{\mathrm{in}}\dot{E}_{x,\mathrm{in}}
-
\sum_{\mathrm{out}}\dot{E}_{x,\mathrm{out}}
-
\dot{W}
=
\dot{E}_{x,\mathrm{dest}}.
\]
对于透平，设备输出功 $\dot{W}_{\mathrm{out}}$ 可由入口和出口火用差扣除火用损失得到：
\[
\dot{W}_{\mathrm{out}}
=
\dot{E}_{x,\mathrm{in}}
-
\dot{E}_{x,\mathrm{out}}
-
\dot{E}_{x,\mathrm{dest}}.
\]
对于压气机或泵，输入功则可表示为
\[
\dot{W}_{\mathrm{in}}
=
\dot{E}_{x,\mathrm{out}}
-
\dot{E}_{x,\mathrm{in}}
+
\dot{E}_{x,\mathrm{dest}}.
\]

由上述表达可知，统一流股对象中的 $h$ 和 $s$ 是连接能量分析与火用分析的关键参数。比焓决定流股携带的能量水平，比熵与参考环境共同决定流股的可用能水平。通过在流股对象中统一定义 $\dot{H}$、$e_x$ 和 $\dot{E}_x$，可以避免在不同设备模型中重复编写能量流和火用流计算公式。

统一流股对象是本文联合循环机组建模方法中的基础环节。它一方面向下连接气体和水蒸气的物性计算方法，另一方面向上支撑各类设备模型和系统性能分析模型。通过该方法，可以使复杂联合循环机组中的多工质、多设备、多状态点问题转化为一组标准化流股对象之间的变换关系，从而提高模型的结构清晰度和工程适用性。

\section{气体工质组成与物性建模方法}

\subsection{燃料气组成表示与归一化处理}

燃料气组成是气体工质建模的基础输入之一。对于燃气轮机系统而言，燃料气并非单一组分，而是由氢气、甲烷、乙烷、一氧化碳、二氧化碳、氮气以及少量高碳烃类共同组成。不同组分的摩尔质量、低位发热量以及热力性质均存在明显差异，因此在进行燃料热值、烟气生成量、混合气焓熵和火用计算之前，需要首先建立统一的燃料气组成表示方法。

CoolProp 物性库是一个开源数据库，提供了多种纯组分和混合物的热力性质计算接口。本文程序利用 CoolProp 查询燃料气中各组分的摩尔质量、低位发热量以及热力性质等参数。为了保证后续计算的一致性，程序预设了一个固定的燃料组分列表，并要求输入数据中的燃料组成必须包含该列表中的所有组分。

气体组分可以表示为各组分的摩尔分数，主要包括
\[
\mathrm{H_2},\ \mathrm{N_2},\ \mathrm{CO_2},\ \mathrm{CH_4},\ \mathrm{CO},\ \mathrm{O_2},\ \mathrm{H_2O},
\mathrm{C_2H_6},\ \mathrm{C_3H_8},\ i\mathrm{C_4H_{10}},\ n\mathrm{C_4H_{10}},
i\mathrm{C_5H_{12}},\ n\mathrm{C_5H_{12}} .
\]
其中，氢气、甲烷、一氧化碳以及 C2--C5 烃类为主要可燃组分，氮气、二氧化碳和氧气等则作为燃料中的惰性组分参与混合物平均性质计算。程序通过固定组分列表保证后续物性计算、燃烧产物计算和低位发热量计算使用一致的组分顺序与组分名称。

由于原始组分数据可能未严格归一化，程序在使用燃料组成前均进行归一化处理。设原始输入中第 \(i\) 种组分的相对含量为 \(y_i\)，则归一化后的摩尔分数 \(x_i\) 为
\[
x_i = \frac{y_i}{\sum_{j=1}^{N} y_j},
\]
其中 \(N\) 为参与建模的气体组分数。归一化处理后，燃料气平均摩尔质量可由各组分摩尔分数加权得到：
\[
M_{\mathrm{mix}} = \sum_{i=1}^{N} x_i M_i,
\]
式中，\(M_i\) 为第 \(i\) 种纯组分的摩尔质量。各纯组分摩尔质量由 CoolProp 物性库查询获得。

\subsection{湿空气组成计算方法}

燃气轮机入口空气是燃烧过程的重要反应物，其组成直接影响燃烧所需氧量、烟气生成量以及压气机和余热锅炉烟气侧的热力性质计算。实际环境空气并非完全干空气，而是含有一定水蒸气的湿空气。环境湿度变化会改变空气中氧气和氮气的有效摩尔分数，并进一步影响燃烧产物中水蒸气含量和烟气平均摩尔质量。因此，在构造燃气轮机入口空气流股时，需要根据环境温度、压力和相对湿度对空气组成进行修正。

首先将干空气视为由氧气、氮气和少量二氧化碳组成的混合气体，其基准摩尔分数取为
\[
x_{\mathrm{O_2,dry}} = 0.2095,\quad
x_{\mathrm{N_2,dry}} = 0.7808,\quad
x_{\mathrm{CO_2,dry}} = 0.0004 .
\]
该组成用于表示进入压气机前的干空气主体成分。由于氩气等微量组分在本文模型中未单独建立物性计算项，其影响在工程计算中近似并入主要惰性组分处理。

湿空气中的水蒸气摩尔分数由环境相对湿度计算。程序中首先调用 CoolProp 根据环境温度 \(T\) 计算水的饱和蒸气压 \(p_{\mathrm{sat}}(T)\)，再由相对湿度 \(RH\) 和环境压力 \(P\) 得到水蒸气分压：
\[
p_{\mathrm{H_2O}} = \varphi p_{\mathrm{sat}}(T),
\]
其中，\(\varphi\) 为相对湿度的小数形式。根据理想气体混合物中摩尔分数等于分压占总压的比例，湿空气中水蒸气摩尔分数可表示为
\[
x_{\mathrm{H_2O}} = \frac{p_{\mathrm{H_2O}}}{P}
=
\frac{\varphi p_{\mathrm{sat}}(T)}{P}.
\]

得到水蒸气摩尔分数后，干空气中各组分的摩尔分数按剩余气相比例进行缩放：
\[
x_{i,\mathrm{wet}} = x_{i,\mathrm{dry}}\left(1-x_{\mathrm{H_2O}}\right),
\quad
i \in \{\mathrm{O_2},\mathrm{N_2},\mathrm{CO_2}\}.
\]
同时令
\[
x_{\mathrm{H_2O,wet}} = x_{\mathrm{H_2O}}.
\]
由此可得湿空气组成：
\[
\boldsymbol{x}_{\mathrm{air}}
=
\left[
x_{\mathrm{O_2,wet}},
x_{\mathrm{N_2,wet}},
x_{\mathrm{CO_2,wet}},
x_{\mathrm{H_2O,wet}}
\right].
\]
随后将该组成转换为统一的气体组成对象，并再次归一化，以消除数值舍入或输入误差对组分总和的影响。

在本文模型中，湿空气组成主要用于两个方面。一方面，入口空气作为燃气轮机压气机工质，需要根据其组成计算焓、熵、定压比热和气体常数等热力性质；另一方面，空气中的氧气、水蒸气和惰性组分会进入燃烧产物计算，决定余热锅炉入口烟气的组成和物性。通过将相对湿度引入空气组成建模，程序能够反映环境条件变化对燃气轮机入口工质和烟气侧热力计算的影响，从而提高不同运行工况下性能分析结果的一致性。

\subsection{燃烧产物与烟气组成计算方法}

燃烧产物组成是燃气轮机排气和余热锅炉烟气侧热力计算的基础。燃料气与空气在燃烧室内反应后，燃料中的可燃组分被氧化生成二氧化碳和水蒸气，空气中的氮气以及燃料中的惰性组分则随烟气进入余热锅炉。烟气组成不仅决定混合气平均摩尔质量，也会影响烟气焓、熵、定压比热、密度以及火用计算结果。因此，在构造烟气流股之前，需要根据燃料组成、空气组成以及二者的质量流量计算燃烧产物组成。

首先利用燃料气和空气的平均摩尔质量，将质量流量换算为摩尔流量：
\[
\dot{n}_{\mathrm{fuel}}
=
\frac{\dot{m}_{\mathrm{fuel}}}{M_{\mathrm{fuel}}},
\qquad
\dot{n}_{\mathrm{air}}
=
\frac{\dot{m}_{\mathrm{air}}}{M_{\mathrm{air}}}.
\]
式中，\(\dot{m}_{\mathrm{fuel}}\) 和 \(\dot{m}_{\mathrm{air}}\) 分别为燃料和空气质量流量，\(M_{\mathrm{fuel}}\) 和 \(M_{\mathrm{air}}\) 分别为燃料气和湿空气平均摩尔质量。

在燃烧反应计算中，程序对燃料中的主要含碳组分和含氢组分分别建立化学计量关系。对于碳元素，甲烷、一氧化碳、乙烷、丙烷、丁烷和戊烷等组分完全燃烧后生成二氧化碳。若第 \(k\) 种燃料组分的摩尔分数为 \(x_k\)，其分子中碳原子数为 \(a_k\)，则由燃料燃烧生成的二氧化碳摩尔流量可表示为
\[
\dot{n}_{\mathrm{CO_2,prod}}
=
\dot{n}_{\mathrm{fuel}}
\sum_k x_k a_k .
\]
对于氢元素，氢气和各烃类燃烧生成水蒸气。若第 \(k\) 种组分完全燃烧生成水蒸气的化学计量系数为 \(b_k\)，则燃烧生成的水蒸气摩尔流量为
\[
\dot{n}_{\mathrm{H_2O,prod}}
=
\dot{n}_{\mathrm{fuel}}
\sum_k x_k b_k .
\]
通过给出各可燃组分对应的碳原子数、水蒸气生成系数和理论需氧系数，将燃烧反应计算转化为统一的组分加权求和。

根据燃烧反应对各组分摩尔流量进行计算后，程序将烟气组分总量归一化，得到烟气摩尔分数组成：
\[
x_{i,\mathrm{flue}}
=
\frac{\dot{n}_{i,\mathrm{flue}}}
{\sum_j \dot{n}_{j,\mathrm{flue}}}.
\]

该方法能够利用运行数据中的燃料流量、空气流量和燃料组分推算余热锅炉入口烟气组成，并为后续烟气焓熵计算、烟气密度计算和余热利用分析提供统一的气体组成输入。

\subsection{混合气体工质的物性计算}

在确定燃料气、湿空气和烟气组成后，需要进一步计算混合气体的焓、熵、定压比热以及气体常数等热力性质，可根据各组分摩尔分数对纯组分物性进行加权，得混合气体的物性。

对于给定温度 \(T\)、压力 \(P\) 和组成 \(\boldsymbol{x}\)，首先对气体组成进行归一化处理，然后遍历所有组分。对于第 \(i\) 种组分，其摩尔分数为 \(x_i\)，程序采用理想气体混合物分压关系计算该组分分压：
\[
p_i = x_i P .
\]
在此基础上，调用纯组分物性计算函数获得该组分在 \(T\)、\(p_i\) 下的摩尔焓 \(h_i^{\mathrm{m}}\)、摩尔熵 \(s_i^{\mathrm{m}}\) 和摩尔定压比热 \(c_{p,i}^{\mathrm{m}}\)。

利用 CoolProp 查询气体纯组分的摩尔质量、摩尔焓、摩尔熵和摩尔定压比热等参数：

\[
h_i^{\mathrm{m}} = h_i^{\mathrm{m}}(T,p_i),\quad
s_i^{\mathrm{m}} = s_i^{\mathrm{m}}(T,p_i),\quad
c_{p,i}^{\mathrm{m}} = c_{p,i}^{\mathrm{m}}(T,p_i).
\]

对于燃料气、空气和烟气中的氢气、氮气、二氧化碳、甲烷、一氧化碳、氧气以及部分烃类组分，通过物性库在给定温度和压力下获取纯组分物性，再完成摩尔分数加权。

对于烟气中的水蒸气，为防止相变取数据库无法处理的问题，采用 Shomate 方程计算其理想气体热力性质。Shomate 方程以温度多项式形式给出水蒸气的摩尔焓、摩尔熵和摩尔定压比热，适用于燃气轮机和余热锅炉烟气侧较宽的温度范围。程序根据温度区间选取不同的 Shomate 系数，并在熵计算中加入压力修正项：
\[
s_{\mathrm{H_2O}}^{\mathrm{m}}
=
s_{\mathrm{H_2O}}^{\mathrm{m},0}(T)
-
R\ln\left(\frac{p_{\mathrm{H_2O}}}{p_{\mathrm{ref}}}\right).
\]
其中，\(p_{\mathrm{ref}}\) 为 Shomate 方程采用的参考压力。该处理使烟气中的水蒸气能够按照理想气体组分参与混合气熵的计算。

混合气体的摩尔焓、摩尔熵和摩尔定压比热按各组分摩尔分数加权得到：
\[
h_{\mathrm{mix}}^{\mathrm{m}}
=
\sum_i x_i h_i^{\mathrm{m}},
\]
\[
s_{\mathrm{mix}}^{\mathrm{m}}
=
\sum_i x_i s_i^{\mathrm{m}},
\]
\[
c_{p,\mathrm{mix}}^{\mathrm{m}}
=
\sum_i x_i c_{p,i}^{\mathrm{m}}.
\]
由于后续设备模型均以质量流量为基础，程序进一步利用混合气平均摩尔质量 \(M_{\mathrm{mix}}\) 将摩尔基准物性转换为质量基准物性：
\[
h_{\mathrm{mix}}
=
\frac{h_{\mathrm{mix}}^{\mathrm{m}}}{M_{\mathrm{mix}}},
\qquad
s_{\mathrm{mix}}
=
\frac{s_{\mathrm{mix}}^{\mathrm{m}}}{M_{\mathrm{mix}}},
\qquad
c_{p,\mathrm{mix}}
=
\frac{c_{p,\mathrm{mix}}^{\mathrm{m}}}{M_{\mathrm{mix}}}.
\]

需要说明的是，由于数据库中气体焓和熵选取的基准并不相同，因此为混合物设置了统一参的考状态。对于每一种组分，程序同时计算其在参考温度 \(T_{\mathrm{ref}}\) 和参考压力 \(P_{\mathrm{ref}}\) 下的物性，并对混合气摩尔焓和摩尔熵进行参考态扣除：
\[
h_{\mathrm{mix}}
=
\frac{
\sum_i x_i h_i^{\mathrm{m}}(T,p_i)
-
\sum_i x_i h_i^{\mathrm{m}}(T_{\mathrm{ref}},p_{i,\mathrm{ref}})
}
{M_{\mathrm{mix}}},
\]
\[
s_{\mathrm{mix}}
=
\frac{
\sum_i x_i s_i^{\mathrm{m}}(T,p_i)
-
\sum_i x_i s_i^{\mathrm{m}}(T_{\mathrm{ref}},p_{i,\mathrm{ref}})
}
{M_{\mathrm{mix}}},
\]
其中，\(p_{i,\mathrm{ref}}=x_iP_{\mathrm{ref}}\)。通过参考态扣除，程序得到的是相对于参考状态的比焓和比熵差值，这与后续气体工质火用计算中的参考环境处理保持一致。

基于上述物性计算结果，当给定温度、压力、质量流量和气体组成时，可以计算得到气体比焓 \(h\)、比熵 \(s\)、定压比热 \(c_p\) 以及气体常数
\[
R_{\mathrm{mix}}=\frac{R}{M_{\mathrm{mix}}},
\]
并将这些参数与流股名称、质量流量和参考环境共同保存为气体流股状态。对于已知压力和焓或已知压力和熵的情况，可以通过二分法迭代计算满足目标焓或目标熵的温度，再通过温度和压力完成状态构造。

综上，本文混合气体物性计算方法以组分摩尔分数为基础，通过纯组分物性计算、理想混合加权、参考态扣除和质量基准转换，获得适用于系统热力分析的焓、熵和定压比热。该方法能够反映燃料组成、空气湿度和燃烧产物组成变化对气体工质热力性质的影响，为后续燃气轮机、余热锅炉及全厂火用分析提供统一的物性基础。

\section{水和水蒸气状态建模方法}

水和水蒸气是联合循环机组汽水系统中的核心工质，其状态变化贯穿余热锅炉、汽轮机、凝汽器、给水泵以及各级减温减压过程。与燃气轮机侧气体工质不同，水和水蒸气在工程运行范围内可能处于压缩液、饱和水、湿蒸汽、干饱和蒸汽和过热蒸汽等不同区域，不能简单采用理想气体模型描述。因此，本文在水和水蒸气工质建模中直接调用成熟的水蒸气物性库，根据输入状态参数计算焓、熵、温度、干度等热力性质。

本文程序采用 Python 的 \verb|iapws| 库计算水和水蒸气物性。该库实现了国际水和水蒸气性质协会提出的 IAPWS-IF97 工业公式，适用于电站热力系统中常见的水和水蒸气状态计算。IAPWS-IF97 将水和水蒸气的热力性质划分为不同区域，并根据压力、温度、焓、熵、干度等输入参数求解状态点。与查表法相比，直接调用 IAPWS-IF97 公式能够提高程序自动化程度，并避免人工插值带来的误差和不便。

对于已知压力和温度的状态点，水和水蒸气状态可表示为压力和温度的函数：
\[
\left(h,s,x\right)=f_{\mathrm{IF97}}(P,T).
\]
由此可得到对应状态的比焓、比熵和干度等参数。这类状态确定方式主要用于过热蒸汽、给水、凝结水以及测点中同时具有压力和温度数据的流股。

对于饱和水或湿蒸汽状态，可采用压力和干度确定状态：
\[
\left(T,h,s\right)=f_{\mathrm{IF97}}(P,x).
\]
当 \(x=0\) 时表示饱和水，当 \(x=1\) 时表示干饱和蒸汽，\(0<x<1\) 时表示湿蒸汽。该方法适合用于汽包饱和水、凝汽器出口凝结水以及汽轮机低压缸排汽等难以直接由温度压力唯一稳定确定的状态点。

对于部分设备，还可采用压力与比焓、压力与比熵两类参数组合反算状态：
\[
\left(T,s,x\right)=f_{\mathrm{IF97}}(P,h),
\qquad
\left(T,h,x\right)=f_{\mathrm{IF97}}(P,s).
\]
这两类状态确定方式在汽轮机等熵膨胀计算、节流过程、混合流股状态重构以及已知能量平衡结果的状态回算中具有较强适用性。

水和水蒸气工质的参考环境同样由 IAPWS97 计算。程序根据参考温度 \(T_0\) 和参考压力 \(P_0\) 构造环境状态，并保存对应的参考比焓 \(h_0\) 与参考比熵 \(s_0\)。这样，水和水蒸气流股不仅能够用于能量流计算，也能够直接用于后续火用流和设备火用损失分析。

综上，本文水和水蒸气工质建模并不重新推导水蒸气状态方程，而是基于 IAPWS-IF97 工业公式库完成基础物性计算，并在程序中封装为统一的流股状态对象。通过 \((P,T)\)、\((P,x)\)、\((P,h)\) 和 \((P,s)\) 四类状态参数组合，模型能够适应汽水系统中不同测点的数据形式，为余热锅炉、汽轮机、凝汽器和给水系统的热力分析提供可靠的状态参数。同时与气体工质的统一流股建模方法一致，为全厂性能分析中的多工质、多设备状态传递提供了基础。

\section{由实测数据构造的流股对象}

为了将实测运行数据直接用于后续设备性能分析，本文根据数据表的单行工况构造全厂流股对象。每一行数据对应一个运行时刻或一个稳态工况，程序读取该行中的温度、压力、质量流量、燃料组分、环境温度和相对湿度等测点信息，分别构造气体侧流股和汽水侧流股。

汽水侧流股主要由两台余热锅炉、汽轮机、凝汽器及相关泵和抽汽流股组成。对于 1 号余热锅炉，构造了 23 个汽水侧流股，包括高压给水泵前后、高压省煤器入口、高压过热器出口、高压减温水、高压减温器后、高压主蒸汽，中压给水泵前后、中压省煤器入口、中压减温水、中压过热器出口，低压省煤器入口和出口、低压汽包给水调节阀出口、低压汽包、低压汽包至低压主蒸汽、低压主蒸汽，以及高压缸排汽、冷再热混合后、再热器出口、再热减温器出口和热再热出口等流股。2 号余热锅炉采用与 1 号余热锅炉相同的流股划分。

此外还构造了汽轮机、凝汽器和热网抽汽相关流股，包括高压缸入口、高压缸出口、中压缸入口、中压缸出口、中压缸出口抽汽后、低压缸入口、低压缸出口、凝汽器出口、凝结水泵出口以及热网抽汽。对于由两台余热锅炉汇合进入汽轮机的流股，程序还根据机组运行状态进行筛选，并按质量流量加权计算混合焓，从而得到进入汽轮机的统一入口流股。

气体侧流股共 10 个，主要覆盖两台燃气轮机及两台余热锅炉烟气侧。具体包括 1 号炉燃料、2 号炉燃料，1 号燃机入口空气、2 号燃机入口空气，1 号燃机压气机出口、2 号燃机压气机出口，1 号余热锅炉入口烟气、2 号余热锅炉入口烟气，以及 1 号余热锅炉出口烟气、2 号余热锅炉出口烟气。

综上，本文在实测数据基础上完成了较完整的全厂流股对象构造工作：两台余热锅炉汽水侧共 46 个流股，全厂汽轮机及凝汽器相关公共汽水侧 10 个流股，燃料、空气和烟气侧 10 个流股，总计 66 个流股对象。通过这一处理，原始运行数据被整理为覆盖燃气轮机、余热锅炉、汽轮机、凝汽器、给水系统和热网抽汽的热力状态集合，为后续设备级能量分析、火用分析和全厂性能评价提供了统一的数据基础。同时也易于维护和扩展，能满足后续研究需求。

\section{本章小结}

本章围绕联合循环机组热力性能分析中的工质建模问题，建立了面向燃料气、空气、烟气、水和水蒸气等多类工质的统一流股建模方法。首先，将不同设备之间传递的工质状态抽象为统一流股对象，使每个流股均包含温度、压力、质量流量、比焓、比熵、组成、能量流和火用流等信息，从而为后续设备模型提供一致的数据接口。

在气体工质方面，本章分别建立了燃料气组成、湿空气组成和燃烧产物烟气组成的计算方法。燃料气模型考虑了氢气、甲烷、一氧化碳以及 C2--C5 烃类等多种组分，并通过归一化处理和平均摩尔质量计算形成统一组成输入；湿空气模型根据环境温度、压力和相对湿度计算水蒸气摩尔分数；烟气模型则基于燃料组成、空气组成和燃烧化学计量关系推算燃烧产物组成。在物性计算方面，程序结合 CoolProp 和 Shomate 方程计算混合气体的焓、熵和定压比热，为燃气轮机和余热锅炉烟气侧分析提供基础。

在水和水蒸气工质方面，本章采用 IAPWS-IF97 工业水蒸气物性公式，通过 \((P,T)\)、\((P,x)\)、\((P,h)\) 和 \((P,s)\) 等状态参数组合确定水和水蒸气状态，避免了手工查表和插值过程。该方法能够适应给水、饱和水、湿蒸汽、过热蒸汽以及混合流股等不同状态点的计算需求，并与统一流股对象结合，用于能量流和火用流计算。

最后，本章结合实测运行数据完成了全厂流股对象构造。程序以单行运行数据为一个工况，最多构造 66 个流股对象，其中汽水侧 56 个、气体侧 10 个，覆盖两台燃气轮机、两台余热锅炉、汽轮机、凝汽器、凝结水泵和热网抽汽等主要系统。通过上述建模工作，原始测点数据被转化为具有明确热力学含义的状态集合，为后续设备级能量分析、火用分析和全厂性能评价奠定了基础。
